\chapter{Einleitung}

Schachcomputer sind faszinierend: Mit erstaunlicher Präzision bewerten sie Stellungen und finden Züge, die für Menschen
manchmal schwer nachvollziehbar sind. Mithilfe der Mathematik haben sie das Spiel in den letzten Jahrzehnten
so stark \enquote{perfektioniert}, sodass sie heute die besten menschlichen Spieler schlagen.

Im modernen Schach sind Schachcomputer daher nicht mehr wegzudenken. Sie werden zur Vorbereitung, zum Analysieren von
Partien und zur Verbesserung des eigenen Spiels eingesetzt.

Als Schachinteressierte wollen wir, Dario und Romeo, uns mit diesem Thema vertieft beschäftigen. Wir möchten verstehen,
wie ein Schachcomputer aufgebaut ist und wie er vorgeht, und anschliessend einen einfachen Schachcomputer entwickeln. \\
Dabei sind wir uns bewusst, dass dieser nicht mit anderen Computern, hinter denen ganze Teams oder Communitys stehen,
mithalten kann. Dennoch finden wir es spannend zu sehen, wie weit wir mit einfachen Mitteln kommen können.

\section{Ausgangslage}
Dabei starten wir nicht von Vorne: Wir durften schon während der Programmierarbeit des CAS-MSED.2501 uns mit der Thematik
auseinandersetzten. Innheralb dieser Arbeit haben wir einen einfachen Schachcomputer programmiert, welcher den MiniMax-Algorithmus verwendet und
die Position anhand den Figurenwerten bewertet. Um gegen den Schachcomputer antreten zu können, haben wir zusätzlich eine
CLI entwickelt. \\
Die Aufgabenstellung befindet sich im Anhang unter \ref{ch:programmierarbeit-fur-cas-msed}. Eine ausführbare JAR Datei
des Schachcomputers, sowie der alte Standpunkt des Codes, kann bei \href{https://github.com/l2c0r3/blobfish/releases/tag/v1.0.0}{GitHub} heruntergeladen werden.

\section{Zielsetzung}
Nun als nächster Schritt, möchten wir uns bei drei Themen weiter vertiefen: Einerseits soll die Zugqualität durch eine
verbesserte \textbf{Bewertungsfunktion} erhöht werden. Dafür möchten wir, neben der Figurenwerte, uns mit weiteren
Heuristiken befassen und hinzufügen. Daneben möchten wir die \textbf{Berechnungszeit}, welcher der Schachcomputer
benötigt, um den nächstbesten Zug berechnen verringern. Und zu guter letzten möchten wir bei der \textbf{Evaluation}
unseren Schachcomputer gegen andere messen. \\
Das ausgefüllte Anmeldeformular befindet sich im Anhang \ref{ch:themeneingabe-transferarbeit}.